% !TEX program = xelatex
% Editable blank template, matching the supplied Word heading structure.
\documentclass[UTF8,fontset=fandol,zihao=-4]{ctexart}
\usepackage{reportstyle}
\newcommand{\slot}[1]{\par{\color{black!55}#1}\par\vspace{6pt}}
\begin{document}
\worktitle{作品名称}
\begin{center}（不要出现任何涉及学校名称等内容）\end{center}
\abstractheading
\slot{800字内}
\section{作品概述}
\subsection{功能与特性}\slot{400字内，可加重要图。}
\subsection{应用领域}\slot{400字内，可加重要图。}
\subsection{主要技术特点}\slot{400字内。}
\subsection{主要性能指标}\slot{200字内，建议用表格。}
\subsection{主要创新点}\slot{200字以内，逐点给出。}
\subsection{设计流程}\slot{200字内，可加重要图。}
\section{系统组成及功能说明}
\slot{阐述具体的设计细节，图文结合。}
\subsection{整体介绍}\slot{给出系统整体框图，标注各子模块，说明模块之间的关系。}
\subsection{硬件系统介绍}
\subsubsection{硬件整体介绍}\slot{介绍硬件整体配置。}
\subsubsection{机械设计介绍}\slot{如果有机械设计，从总体到局部给出组件设计图，可使用 CAD 截图或手绘图。}
\subsubsection{电路各模块介绍}\slot{逐级说明各模块电路，标记关键输入输出信号线，可附原理图或 PCB 图。}
\subsection{软件系统介绍}
\subsubsection{软件整体介绍}\slot{含 PC 端或云端，结合关键图片。}
\subsubsection{软件各模块介绍}\slot{从顶层到底层说明模块与函数流程，列出关键输入输出变量。}
\section{完成情况及性能参数}
\slot{阐述最终成果，图文结合，以实物照片为主。}
\subsection{整体介绍}\slot{提供整个系统实物正面和斜45度全局照片。}
\Needspace{10\baselineskip}
\subsection{工程成果}
\subsubsection{机械成果}\slot{实物照片。}
\subsubsection{电路成果}\slot{实物照片。}
\subsubsection{软件成果}\slot{界面照片。}
\subsection{特性成果}\slot{展示功能与性能等量化指标，可附仪器测试或现场照片。}
\section{总结}
\subsection{可扩展之处}\slot{300字内。}
\subsection{心得体会}\slot{1000字内，可包括研发和制作细节。}
\section{参考文献}\slot{按照标准格式，20篇以内。}
\end{document}
